\scanpage{122}
\mspage{109}
et \deleted{\unsure{on}} a égalité si $I$ est \unsure{noethérien}, i.e. $I^\circ$ est
artinien.

\uline{Dém}. La démonstration de $\deleted{\unsure{\Phi_I}}\Phi(I)\subset\Phi_I$
est celle qu'on vient de donner. Il reste
à prouver que si $I$ est \deleted{artinien}\inserted{noethérien}, on a
$\Phi_I\subset\Phi(I)$. Etant donné $\unsure{i|u}\colon A\longrightarrow B\unsure{=}$
dans $\Phi_I$, il faut prouver que $u\in\Phi(I)$, i.e.
que pour un diagramme comm. dans $M(I)$,

\diagram{\((6.3.2)\)\quad
\deleted{\begin{tikzcd}[ampersand replacement=\&]
A\arrow[r,"i"]\arrow[d] \& B\arrow[d]\arrow[dl,dashed]\\
X\arrow[r,"p"] \& Y
\end{tikzcd}}\qquad
\begin{tikzcd}[ampersand replacement=\&]
A\arrow[r,"\alpha"]\arrow[d,"i"'] \& X\arrow[d,"p"]\\
B\arrow[r,"\beta"']\arrow[ur,dashed,"h"] \& Y
\end{tikzcd}}

avec $p\in\Psi\langle I\rangle$ i.e. $p(i)\colon X(i)\unsure{\text{---}|\longrightarrow}Y(i)$
dans $\Psi$ pour \unsure{tt} $i$, $\exists h\colon B\longrightarrow X$ laissant le
diagramme commutatif. On va revenir
aux notations duales, en travaillant dans
$M(I^\circ)$, avec $I$ artinien. On considère
l'\unsure{ens} \inserted{$\unsure{\mathcal H|H}$} des couples $(U,h_U)$ d'un ouvert $U$
de $I$ et d'un $h_U\colon B|U\longrightarrow X|U$,
rendant commutatif la restriction à $U$
du \unsure{car.|carr.} On ordonne $\mathcal H$ par prolongement,
on voit que tte \deleted{partie} tot. ordonnée \inserted{$(U_\alpha,h_{U_\alpha})$} dans
$\mathcal H$ a une borne supérieure, savoir
$U=\bigcup_\alpha U_\alpha$, $h_U$ défini \deleted{par les} $h_{U_\alpha}$ sur les $U_\alpha$.
% archive sheet number: 121

